\subsection{不稳定轨迹的分层}\label{src:9.3.3-subsection}
凝聚层的模叠 $\Coh_{r,d}$ 有以 Harder--Narasimhan 多边形为指标的局部闭子叠分层. 此分层称为\newterm{Harder--Narasimhan 分层}{Harder--Narasimhan stratification} 或\newterm{Shatz 分层}{Shatz stratification}. 它最早在 \cite{Sha77} 中证明, 与本章大多数其他结果 (例如 Harder--Narasimhan 滤过及 Langton 定理) 一样, 可推广到任意维射影簇上的无挠层.

回顾对 $C$ 上凝聚层 $F$, Harder--Narasimhan (HN) 多边形 $\operatorname{HNP}(F)\subseteq\mathbb C$ 是当 $E\subseteq F$ 遍历所有不稳定化子层 (即 $\mu(E)\geq\mu(F)$, 等价地, $\arg Z(E)\geq\arg Z(F)$) 时, 点 $Z(E)=-\deg(E)+\operatorname{rk}(E)i$ 的凸包; 见图~\bookref{9.1.25}. 若 $0=F_0\subseteq F_1\subseteq\cdots\subseteq F_l=F$ 是 Harder--Narasimhan 滤过 (\bookref{9.1.21}), 则 $Z(F_i)$ 是 $\operatorname{HNP}(F)$ 的极端点 (推论~\bookref{9.1.26}). HN 多边形由因子 $F_i/F_{i-1}$ 的秩次数对序列确定. 我们使用由包含关系给出的 HN 多边形的偏序. 注意若 $F$ 半稳定, 则 $\operatorname{HNP}(F)$ 为从 $0$ 到 $Z(F)$ 的直线; 若 $F$ 为挠层, 则 $\operatorname{HNP}(F)$ 为顶点 $(0,0),(-\deg F,0),Z(F)$ 构成的多边形.
\begin{theorem}[Harder--Narasimhan/Shatz 分层]\label{src:9.3.15}
设 $C$ 是代数闭域 $\Bbbk$ 上亏格 $g\geq2$ 的光滑连通射影曲线. 存在分层
\[
\Coh_{r,d}=\coprod_{\text{HN 多边形 }\mathcal P}\Coh_{r,d}^{\mathcal P},
\]
其中 $\Coh_{r,d}^{\mathcal P}\subseteq\Coh_{r,d}$ 为局部闭子叠, 对 $C$ 上秩为 $r$, 次数为 $d$ 的向量丛 $E$, 有
\[
E\in\Coh_{r,d}^{\mathcal P}\iff\operatorname{HNP}(E)=\mathcal P.
\]
此外, 对每个 HN 多边形 $\mathcal P$:
\begin{enumerate}[(1)]
\item 设 $S$ 是 $\Bbbk$-概形, $\mathcal F$ 是 $C_S$ 上秩为 $r$, 次数为 $d$ 的向量丛族.
\begin{enumerate}[(a)]
\item 分类 $\mathcal F$ 的映射 $S\to\Coh_{r,d}$ 经过 $\Coh_{r,d}^{\mathcal P}$ 分解, 当且仅当存在向量丛滤过
\[
0=\mathcal F_0\subseteq\mathcal F_1\subseteq\cdots\subseteq\mathcal F_l=\mathcal F,
\]
其因子 $\mathcal F_i/\mathcal F_{i-1}$ 为向量丛, 对每个 $s\in S$, 有 $\operatorname{HNP}(\mathcal F_s)=\mathcal P$, 且 $0=\mathcal F_{0,s}\subseteq\mathcal F_{1,s}\subseteq\cdots\subseteq\mathcal F_{l,s}=\mathcal F_s$ 是 $\mathcal F_s$ 的 Harder--Narasimhan 滤过.
\item 若 $s\leadsto s_0$ 为特殊化, 则 $\operatorname{HNP}(\mathcal F_s)\subseteq\operatorname{HNP}(\mathcal F_{s_0})$; 换言之, $\Coh_{r,d}^{\mathcal P}$ 的闭包包含于 $\coprod_{\mathcal P'\supseteq\mathcal P}\Coh_{r,d}^{\mathcal P}$.
\end{enumerate}
\item 并集
\[
\Coh_{r,d}^{\leq\mathcal P}:=\bigcup_{\mathcal P'\subseteq\mathcal P}\Coh_{r,d}^{\mathcal P'}
\]
为 $\Coh_{r,d}$ 的开子叠, 而 $\Coh_{r,d}^{\mathcal P}\subseteq\Coh_{r,d}^{\leq\mathcal P}$ 为闭子叠.
\item 以 $e_1,\ldots,e_l\in\mathbb C$ 表示按虚部递增排列的 $\mathcal P$ 的极端点, 以 $P_1,\ldots,P_l$ 表示满足 $Z(G_i)=e_i-e_{i-1}$ 的向量丛 $G_i$ 的 Hilbert 多项式, 则有等同
\[
\Coh_{r,d}^{\mathcal P}\cong
\operatorname{Flag}^{P_1,\ldots,P_l}(\mathcal F_{\mathrm{univ}}/C\times\Coh_{r,d}^{\leq\mathcal P}/\Coh_{r,d}^{\leq\mathcal P}),
\]
其中右端为 $\Coh_{r,d}^{\leq\mathcal P}$ 上的 Flag 概形 (习题~\bookref{2.4.7}), 参数化万有丛 $\mathcal F_{\mathrm{univ}}$ 的滤过 $\mathcal F_\bullet$, 第 $i$ 个因子 $\mathcal F_i/\mathcal F_{i-1}$ 的 Hilbert 多项式为 $P_i$.
\end{enumerate}
\end{theorem}
\begin{remark}\label{src:9.3.16}
当 $\mathcal P$ 为从 $0$ 到 $Z(F)$ 的直线时, $\Coh_{r,d}^{\mathcal P}=\cBun_{r,d}^{\mathrm{ss}}$ 是半稳定轨迹.
\end{remark}
\begin{proof}
先证明对 $\Bbbk$-概形 $S$ 以及 $C_S$ 上在 $S$ 上平坦的秩 $r$, 次数 $d$ 凝聚层 $\mathcal F$, 存在分层
\[
S=\coprod_{\mathcal P}S_{\mathcal P},
\]
分为局部闭子概形 $S_{\mathcal P}\subseteq S$, 用 $S_{\mathcal P}$ 替换 $\Coh_{r,d}^{\mathcal P}$ 后, 性质 (1)--(3) 成立; 当 $S$ 为 Noether 时分层有限. 由极限方法 (\bookref{B.3}), 可假设 $S$ 为 Noether.

纤维 $\mathcal F_s$ 的 Harder--Narasimhan 多边形只有有限种可能的 $\mathcal P$, 因为纤维子层的次数上有界 (习题~\bookref{9.1.28}). 其次, 对每个这样的多边形 $\mathcal P$, 其极端顶点有 Hilbert 多项式 $P_1,\ldots,P_l$ (沿用 (3) 的记号), 存在 $S$ 上固有的 Flag 概形 $\operatorname{Flag}^{P_1,\ldots,P_l}(\mathcal F/C_S/S)$. 将其闭像记为
\[
S_{\geq\mathcal P}:=\im(\operatorname{Flag}^{P_1,\ldots,P_l}(\mathcal F/C_S/S)\to S).
\]
点 $s\in S$ 位于 $S_{\geq\mathcal P}$, 当且仅当 $\mathcal F_s$ 有伴随多边形为 $\mathcal P$ 的滤过. 则
\[
S_{\mathcal P}:=S_{\geq\mathcal P}\setminus\bigcup_{\mathcal P'\supsetneq\mathcal P}S_{\geq\mathcal P'}
\]
为 $S_{\geq\mathcal P}$ 的开子叠. 因为 Harder--Narasimhan 多边形包含每个不稳定化子层 (推论~\bookref{9.1.26}), 点 $s\in S$ 位于 $S_{\mathcal P}$ 当且仅当 $\operatorname{HNP}(\mathcal F_s)=\mathcal P$. 除 $S_{\mathcal P}$ 的函子性外, 这给出 (1)(a).

为证明 (1)(b), 设 $\Spec R\to S$ 为表示特殊化的来自 DVR 的映射. 由于 Harder--Narasimhan 滤过在域扩张下保持不变, 只需证明 $\operatorname{HNP}(\mathcal F_K)\subseteq\operatorname{HNP}(\mathcal F_\kappa)$, 其中 $K,\kappa$ 为 $R$ 的分式域与剩余域. 由 Quot 概形的赋值判据 (命题~\bookref{2.4.2}), $\mathcal F_K$ 的 Harder--Narasimhan 滤过延拓为 $\mathcal F_R$ 的滤过, 再限制为 $\mathcal F_\kappa$ 的滤过. 因为 $\mathcal F_\kappa$ 的此滤过的每一项都位于 $\operatorname{HNP}(\mathcal F_\kappa)$ 中, 有 $\operatorname{HNP}(\mathcal F_K)\subseteq\operatorname{HNP}(\mathcal F_\kappa)$. 由于 $S_{\leq\mathcal P}:=\bigcup_{\mathcal P'\subseteq\mathcal P}S_{\mathcal P'}$ 可构造且在一般化下稳定, 它是开的, 从而建立 (2).

剩下需证明 (2) 以及概形结构 $S_{\mathcal P}\subseteq S$ 对 $S$ 函子性; 只需证明
\[
S_{\mathcal P}\cong\operatorname{Flag}^{P_1,\ldots,P_l}(\mathcal F/C\times S_{\leq\mathcal P}/S_{\leq\mathcal P}).
\]
换言之, 需证明固有态射
\[
\operatorname{Flag}^{P_1,\ldots,P_l}(\mathcal F/C\times S_{\leq\mathcal P}/S_{\leq\mathcal P})\to S_{\leq\mathcal P}
\]
是同构地映到 $S_{\mathcal P}$ 的闭浸入. 为此, 只需说明它在切空间上单射, 即相对切空间为零. 可将上述映射分解为复合
\[
\begin{aligned}
&\operatorname{Flag}^{P_1,\ldots,P_l}(\mathcal F/C\times S_{\leq\mathcal P}/S_{\leq\mathcal P})\\
&\quad\to\operatorname{Flag}^{P_2,\ldots,P_l}(\mathcal F/C\times S_{\leq\mathcal P}/S_{\leq\mathcal P})\to\cdots\\
&\quad\to\Quot^{P_l}(\mathcal F/C\times S_{\leq\mathcal P}/S_{\leq\mathcal P})\to S_{\leq\mathcal P},
\end{aligned}
\]
其中每个态射由相对 Quot 概形给出 (见习题~\bookref{2.4.7}). 分析上述各映射的相对切空间后, 可将单射性断言归约为: 若 $0=\mathcal F_{s,0}\subseteq\mathcal F_{s,1}\subseteq\cdots\subseteq\mathcal F_{s,l}=\mathcal F_s$ 是 $s\in S_{\mathcal P}$ 上纤维 $\mathcal F_s$ 的 Harder--Narasimhan 滤过, $q:\mathcal F_{s,i}\to\mathcal F_{s,i}/\mathcal F_{s,i-1}$ 为第 $i$ 个商, 则 $T_q\Quot^{P_i}(\mathcal F_{s,i}/C_s/\kappa(s))=0$. 由习题~\bookref{2.5.6}, 此切空间与 $\Hom_{\mathcal O_{C_s}}(\mathcal F_{s,i},\mathcal F_{s,i}/\mathcal F_{s,i-1})$ 等同. 因为 $\mathcal F_{s,i}$ 的 Harder--Narasimhan 滤过的每个因子的斜率都严格大于半稳定层 $\mathcal F_{s,i}/\mathcal F_{s,i-1}$ 的斜率, 不存在非平凡同态 (习题~\bookref{9.1.31}). 另见 \cite{Sha77,Nit11}.
\end{proof}
\begin{remark}[相对 Harder--Narasimhan 滤过]\label{src:9.3.17}
由此定理, 对 $C_S$ 上向量丛 $\mathcal F$, 其限制 $\mathcal F_{S_{\mathcal P}}$ 有向量丛滤过 $0=\mathcal F_0\subseteq\mathcal F_1\subseteq\cdots\subseteq\mathcal F_l=\mathcal F_{S_{\mathcal P}}$, 因子 $\mathcal F_i/\mathcal F_{i-1}$ 为向量丛. 对每个 $s\in S_{\mathcal P}$, 此滤过限制为 $C_s$ 上 $\mathcal F_s$ 的 Harder--Narasimhan 滤过, 因而称为\newterm{相对 Harder--Narasimhan 滤过}{relative Harder--Narasimhan filtration}. 取 $\mathcal P$ 为纤维 $\mathcal F_s$ 的 Harder--Narasimhan 多边形中的极小者, 得到开子概形 $U:=S_{\mathcal P}\subseteq S$, 在其上存在相对 Harder--Narasimhan 滤过. 一般而言, 此滤过未必能在 $S$ 上延拓, 但作一次爆破后可以延拓 \cite[定理 2.3.2]{HL10}.
\end{remark}
\begin{remark}[与 HKKN 分层的关系]\label{src:9.3.18}
$\Coh_{r,d}$ 的 Harder--Narasimhan/Shatz 分层与构造模空间 $M_{r,d}^{\mathrm{ps}}$ 的 GIT 设置中的 HKKN 分层一致 \cite{Hos18}.
\end{remark}
\begin{remark}[\texorpdfstring{$\Theta$}{Θ}-分层]\label{src:9.3.19}
$\Coh_{r,d}$ 的 Harder--Narasimhan/Shatz 分层是注记~\bookref{8.7.8} 中定义的 $\Theta$-分层的例子. 对每个 HN 多边形 $\mathcal P$, $\Coh_{r,d}^{\leq\mathcal P}:=\bigcup_{\mathcal P'\leq\mathcal P}\Coh_{r,d}^{\mathcal P'}\subseteq\Coh_{r,d}$ 为开子叠, 并有等同
\[
\begin{aligned}
&\underline{\operatorname{Mor}}([\mathbb A^1/\mathbb G_m],\Coh_{r,d}^{\leq\mathcal P})\\
&\qquad\cong\coprod_{\mathcal P'}\operatorname{Flag}^{\mathcal P'}(\mathcal F_{\mathrm{univ}}/C\times\Coh_{r,d}^{\leq\mathcal P}/\Coh_{r,d}^{\leq\mathcal P}),
\end{aligned}
\]
其中 $\operatorname{Flag}^{\mathcal P'}$ 是上述 (3) 中 $\operatorname{Flag}^{P_1',\ldots,P_l'}$ 的简写; 此事实强化了来自 $[\mathbb A^1/\mathbb G_m]$ 的映射与滤过之间的对应 (命题~\bookref{7.9.6}). 取 $\mathcal P'=\mathcal P$, Flag 概形 $\operatorname{Flag}^{\mathcal P}(\mathcal F_{\mathrm{univ}}/C\times\Coh_{r,d}^{\leq\mathcal P}/\Coh_{r,d}^{\leq\mathcal P})\subseteq\underline{\operatorname{Mor}}([\mathbb A^1/\mathbb G_m],\Coh_{r,d}^{\leq\mathcal P})$ 为连通分支之并, 且由 (3) 定义一个 $\Theta$-层, 即诱导映射 $\operatorname{Flag}^{\mathcal P}(\mathcal F_{\mathrm{univ}}/C\times\Coh_{r,d}^{\leq\mathcal P}/\Coh_{r,d}^{\leq\mathcal P})\to\Coh_{r,d}^{\leq\mathcal P}$ 是同构地映到 $\Coh_{r,d}^{\mathcal P}$ 的闭浸入.

一般地, 若代数叠 $\mathcal X$ 满足固有性的赋值判据的存在性部分, 且 $\mathcal X=\coprod_\lambda\mathcal S_\lambda$ 是 $\Theta$-分层, 则半稳定轨迹 $\mathcal X^{\mathrm{ss}}:=\mathcal X_{\leq0}=\mathcal S_0$ 也满足固有性的赋值判据的存在性部分 \cite[定理 6.5]{AHLH23}. 因而 Langton 定理 (\bookref{9.3.1}) 也可由 $\Coh_{r,d}$ 的 Harder--Narasimhan/Shatz 分层 (\bookref{9.3.15}) 推出.
\end{remark}
